\section{Introducción}
\label{sec:introduccion}

En el campo de la Inteligencia Artificial, las redes neuronales artificiales
han sido históricamente una de las principales herramientas para el modelado
de funciones complejas y el aprendizaje a partir de datos. En cursos
introductorios de la disciplina, se estudian arquitecturas clásicas como el
perceptrón, el perceptrón multicapa (MLP), las redes de Hopfield y los mapas
autoorganizados (SOM), las cuales permiten comprender los fundamentos del
aprendizaje supervisado, no supervisado y asociativo.

Sin embargo, estas arquitecturas presentan limitaciones estructurales cuando
se enfrentan a problemas que requieren el modelado de dependencias contextuales
complejas, como es el caso del procesamiento del lenguaje natural (NLP, por sus
siglas en inglés). El lenguaje humano es inherentemente secuencial, ambiguo y
dependiente del contexto, características que dificultan su tratamiento mediante
modelos diseñados para entradas de tamaño fijo o relaciones puramente locales.

En este contexto, la arquitectura Transformer, introducida por
\textcite{vaswani2017attention}, representa un cambio de paradigma en el diseño
de redes neuronales para el procesamiento de secuencias, al basarse en el
mecanismo de atención en lugar de recurrencia o convolución. Su impacto fue tal
que dio origen a la familia de modelos de lenguaje de gran escala que hoy
sustentan aplicaciones como traducción automática, generación de texto y
asistentes conversacionales~\cite{rothman2022transformers}.

Este trabajo tiene como objetivo analizar los principios teóricos que
fundamentan la arquitectura Transformer, contrastándolos con los modelos
clásicos estudiados en la cátedra, y desarrollar un prototipo mínimo funcional
que permita evidenciar dichas diferencias de manera concreta. El prototipo
consiste en un modelo de lenguaje autoregresivo a nivel carácter, construido
desde cero en PyTorch, sin recurrir a implementaciones de alto nivel de la
librería.

\subsection*{Alcance y nivel de tratamiento}

Las arquitecturas y conceptos presentados en este trabajo se describen con
un nivel de detalle acorde al carácter introductorio de la cátedra. En
particular, los modelos clásicos (perceptrón, MLP, Hopfield, SOM) se
retoman como punto de contraste conceptual, sin pretender agotar su
formalización matemática completa. Del mismo modo, la descripción del
Transformer y su implementación apuntan a la comprensión de los principios
fundamentales, y no a reproducir el estado del arte en escala o rendimiento.
El objetivo es siempre la claridad conceptual por sobre la exhaustividad
técnica.

\subsection*{Organización del documento}

El documento se organiza de la siguiente manera:
la Sección~\ref{sec:redes_clasicas} presenta las arquitecturas clásicas y sus
limitaciones estructurales;
la Sección~\ref{sec:problema_nlp} caracteriza el problema del procesamiento del
lenguaje natural;
la Sección~\ref{sec:mecanismo_atencion} introduce el mecanismo de atención como
respuesta a dichas limitaciones;
la Sección~\ref{sec:arquitectura_transformer} describe en detalle la arquitectura
Transformer;
la Sección~\ref{sec:transformer_minimo} define formalmente el modelo implementado;
la Sección~\ref{sec:conclusiones} presenta las conclusiones del trabajo; y el
Anexo~\ref{sec:implementacion} incluye la implementación comentada.
